\documentclass[12pt]{article}
\usepackage[utf8]{inputenc}
\usepackage[T1]{fontenc}
\usepackage{hyperref}
\title{Tulehu Language Resources for K-12 Curriculum Development}
\author{Compiled by Claude Sonnet 5 (claude-sonnet-5) for Andr\'as Kornai}
\date{}
\begin{document}
\maketitle

\section{Introduction}
Tulehu (ISO 639-3 tlu; Glottocode tule1244), also called Sou Amana Teru, is an Austronesian language of the Ambonic group, spoken on the northeast coast of Ambon Island, Maluku, Indonesia (Tulehu, Tengah-Tengah, Tial; Liang is also cited). Speaker figures vary: 19,000 (Wikipedia, 1987), 18,800 (Verbix), 10,000 (ASJP). I found no school materials or detailed grammar.

\section{Primary Educational Resources}
None found.

\section{Linguistic and Reference Materials}
\begin{itemize}
\item Wikipedia, Tulehu language: \url{https://en.wikipedia.org/wiki/Tulehu_language}
\item Verbix Tulehu page: \url{https://docs.verbix.com/Languages/Tulehu}
\item Glottolog languoid page: \url{https://glottolog.org/resource/languoid/id/tule1244.bigmap.html}
\item Austronesian Basic Vocabulary Database, Sou Amana Teru: \url{https://abvd.eva.mpg.de/austronesian/language.php?id=379}
\end{itemize}

\section{Religious and Community Resources}
Joshua Project has a page on the Northeast Ambonese Tulehu. It says the people speak Ambonese Malay, so use it with caution: \url{https://joshuaproject.net/people_groups/15625}

\section{Digital and Audio Resources}
\begin{itemize}
\item ASJP wordlist, Tulehu: \url{https://asjp.clld.org/languages/TULEHU}
\end{itemize}

\section{Curriculum Development Context}
Only brief descriptions and wordlists were found. A curriculum would need new documentation, and the language's relationship to Ambonese Malay in daily use should be examined.

\end{document}